\subsection{All smooth obstacles}
\label{subsec:all-obstacle-topology}

The preceding failure comes from choosing one obstacle. Requiring convergence for every smooth obstacle makes the collection stable under a change of background. The same dyadic approximation still determines its analytic closure, now by testing envelope equality at every obstacle.

Let $\tau^{\mathrm{all}}_\omega$ be the uniform topology generated by all the uniformities $\tau_{\omega,F}$, with $F\in C^\infty(X,\mathbb R)$. Thus an entourage uses finitely many obstacle-envelope distances, together with finitely many of the ideal relations and divisorial distances already defined.

\begin{proposition}\label{prop:all-obstacle-uniformity}\index{Singularity type!all smooth obstacles}
A countable uniformly dense family $\{F_k\}_{k\geq1}\subset C^\infty(X,\mathbb R)$ suffices for the obstacle coordinates of $\tau^{\mathrm{all}}_\omega$. For $\chi$ smooth with $\omega'=\omega+\ddc\chi>0$, the map $[u]\mapsto[u-\chi]$ is a uniform isomorphism between $\tau^{\mathrm{all}}_\omega$ and $\tau^{\mathrm{all}}_{\omega'}$. Positive scaling gives the corresponding uniform isomorphism to $\tau^{\mathrm{all}}_{s\omega}$. In positive dimension this topology is not $T_0$ on ordinary classes.
\end{proposition}
\begin{proof}
For either envelope operator, the obstacle comparison gives
\[
 \|P_{\omega,F}[u]-P_{\omega,G}[u]\|_{L^1(\mathrm dV)}
 \leq\|F-G\|_\infty.
\]
The difference between the two envelope distances for a pair $[u],[v]$ is therefore at most $2\|F-G\|_\infty$, uniformly in the pair. Smooth functions contain a countable dense subset for the uniform norm: approximate each member of a countable dense subset of $C(X)$ by smooth functions with errors $1/m$, $m\geq1$. Approximating each of the finitely many obstacles in an entourage and then shrinking its radii proves equality of the uniformities. The divisorial models are still allowed to range over every fixed compact modification.

The exact transport in \cref{prop:ordinary-topology-background-transport} reindexes the entire smooth-obstacle family by $F\mapsto F-\chi$, or by $F\mapsto sF$. These are bijections of that family, proving both isomorphisms.

For nonseparation, use the unbounded $u=1-\sqrt{1-w}$ in \cref{prop:ordinary-topology-nonseparation}. It has trivial ideals and zero divisorial numbers, and $\delta w\preceq u$ for $0<\delta\leq1$. Fix $F$ and put $q=P_\omega(F)$. Since $\omega>0$, the constant $\min_XF$ is a competitor, so $q$ is bounded and $q\leq F$. The function
\[
 h_\delta=(1-\delta)q+\delta w-\delta\|F\|_\infty
\]
is $\omega$-plurisubharmonic, lies below $F$, and satisfies $h_\delta\preceq u$. Letting $\delta\downarrow0$ off the pole of $w$ gives $P_{\omega,F}[u]=q=P_{\omega,F}[0]$. Trivial ideals give the same equality for the ideal envelopes. Every generating entourage therefore has identical sections at $[u]$ and $[0]$, although these ordinary classes differ.
\end{proof}

\begin{theorem}\label{thm:all-obstacle-analytic-closure}\index{Analytic singularity!closure for all obstacles}
The closure and sequential closure of $\mathcal A_\omega$ in $\tau^{\mathrm{all}}_\omega$ both equal
\[
 \mathcal G^{\mathrm{all}}_\omega
 =\{[u]:P_{\omega,F}[u]=P_{\omega,F}[u]_{\mathcal I}
                      \text{ for every smooth }F\}.
\]
It suffices to test equality on the countable dense family above. The dyadic sequence $b_j$ of \cref{cor:fixed-hermitian-analytic-density} converges to $[u]$ in this topology if and only if $[u]\in\mathcal G^{\mathrm{all}}_\omega$.
\end{theorem}
\begin{proof}
Fix a smooth obstacle $F$. The increasing-maxima proof of \cref{thm:hermitian-i-projection} applies to the family below $F$: a translate of $u$ is a lower competitor, and the fixed upper bound supplies the local convergence needed for increasing-limit stability of ideals. Thus $P_{\omega,F}[u]_{\mathcal I}\sim_{\mathcal I}u$. For an analytic $a$, \cref{lem:analytic-ideal-dominance} makes this ideal envelope an ordinary competitor, so $P_{\omega,F}[a]=P_{\omega,F}[a]_{\mathcal I}$.

Write $p=P_{\omega,F}[u]_{\mathcal I}$ and $q_j=P_{\omega,F}[b_j]_{\mathcal I}$. The dyadic construction gives $u\preceq_{\mathcal I}b_j$, decreasing ordinary types, and convergence of every local divisorial Lelong number to that of $u$. Ideal order, proved by taking maxima with a translate of $b_j$ below $F$, gives $p\leq q_j\leq F$. Analytic envelope equality and the decreasing types make $q_j$ decrease. Its limit $q$ is genuine and $\omega$-plurisubharmonic, with $p\leq q\leq F$; the same bounds give convergence in $L^1$.

Pullback convergence and Lelong semicontinuity, as in \cref{prop:ordered-ideal-envelope-convergence}, give $\nu_E(q)\geq\nu_E(u)$ for every local divisor $E$. The inequality $p\leq q$ gives the reverse bound. The local ideal criterion yields $q\sim_{\mathcal I}u$. Since $q\leq F$, it is admissible for $p$, proving
\[
 P_{\omega,F}[b_j]=P_{\omega,F}[b_j]_{\mathcal I}
       \searrow P_{\omega,F}[u]_{\mathcal I}.
\]
The controlled ideal relations and all fixed-model divisorial distances converge by the proof of \cref{thm:ordinary-topology-analytic-closure}; they do not depend on $F$. This one dyadic sequence consequently converges in $\tau^{\mathrm{all}}_\omega$ exactly under the stated envelope equalities.

For each $F$, the equality of the two $L^1$ envelope coordinates is closed. Analytic classes belong to every such equality set, so their net closure lies in the displayed intersection. The constructed sequence proves the converse. Finally the uniform obstacle bound in the preceding proof extends equality from a dense family to every smooth obstacle.
\end{proof}

The inclusion $\mathcal G^{\mathrm{all}}_\omega\subseteq\mathcal G_\omega$ follows by taking $F=0$. On curves the two sets coincide, as proved in \cref{prop:curve-all-obstacle-topology}. The theorem does not identify them in higher dimension. Although $\tau^{\mathrm{all}}_\omega$ removes the fixed-obstacle gauge dependence, it still has incomplete curve families and remains strictly weaker than polar total variation on ordinary curve models, by \cref{prop:curve-all-obstacle-incomplete}.
